\subsection{Плотность}
%%%%%%%%%%%%%%%%%%%%%%%%%%%%%%%%%%%%%%%%%%%%%%%%%%%%
%%%%%%%%%%%%%%%%%%%%%%%%%%%%%%%%%%%%%%%%%%%%%%%%%%%%
%%%%%%%%%%%%%%%%%%%%%%%%%%%%%%%%%%%%%%%%%%%%%%%%%%%%
\z{Высота Эйфелевой башни $H = 324\,\text{м}$, масса $M = 10000\,\text{т}$. Чему будет равна масса статуэтки в виде Эйфелевой башни, если она является её точной копией, уменьшенной в $k = 400$ раз? Ответ выразить в граммах.}
%
{$m = 156.25\,\text{г}$}

%%%%%%%%%%%%%%%%%%%%%%%%%%%%%%%%%%%%%%%%%%%%%%%%%%%%
\z[{\begin{wrapfigure}[5]{r}{0.25\textwidth}%
		\vspace{-25pt}%
		\begin{flushright}%
 			\includegraphics[width=0.25\textwidth]{1.2.2}%
		\end{flushright}%
	\end{wrapfigure}%
	}]%
{Три детали, изготовленные из разных материалов, склеили так, что получился составной шарик. Объёмы, плотности и масса для некоторых деталей указаны на рисунке, где $M$ и $V$ -- масса и объём составного шарика соответственно. Какова средняя плотность шарика?}
%
{$\rho_\text{ср}=\frac{3}{8}(2\rho_1+\rho_2)$}

%%%%%%%%%%%%%%%%%%%%%%%%%%%%%%%%%%%%%%%%%%%%%%%%%%%%
\z{Высота уровня в цилиндрической бочке $h = 1\,\text{м}$. В бочку засыпают маленькие шарики, и вода покрывает их с верхом. При этом плотность смеси $4070\,\text{кг/м}^3$. Высота воды в бочке теперь $1.5\,\text{м}$. Плотность воды $\rho_{\text{в}} = 1000\,\text{кг/м}^3$. Найти плотность шариков.}
%
{$\rho_{\text{ш}} = 10210\,\text{кг}/\text{м}^3$}

%%%%%%%%%%%%%%%%%%%%%%%%%%%%%%%%%%%%%%%%%%%%%%%%%%%%
\z[{\begin{wrapfigure}[5]{r}{0.3\textwidth}%
		\vspace{-10pt}%
		\begin{flushright}%
 			\includegraphics[width=0.3\textwidth]{1.2.4}%
		\end{flushright}%
	\end{wrapfigure}%
	}]%
{Если в кубе массой $m_0=1.6 \,\text{кг}$ сделать лунку в форме кубика, то он будет иметь массу $m_1=1.2 \,\text{кг}$. А если эту лунку заполнить водой, то масса куба будет равна $m_2=1.7 \,\text{кг}$. Определите плотность материала $\rho_x$, из которого изготовлен куб.}
%
{$\rho = 800\,\text{кг/м}^3$}

%%%%%%%%%%%%%%%%%%%%%%%%%%%%%%%%%%%%%%%%%%%%%%%%%%%%
\z{Кубики сахара-рафинада плотно упакованы в коробку, на которой написано: «масса сахара $500\,\text{г}$, $168$ штук». Длина самого большого ребра коробки $98\,\text{мм}$. Вдоль самого короткого ребра коробки укладывается ровно $4$ кусочка сахара. Чему равна плотность сахара-рафинада?}
%
{$\rho = 1080\,\text{кг}/\text{м}^3$}

%%%%%%%%%%%%%%%%%%%%%%%%%%%%%%%%%%%%%%%%%%%%%%%%%%%%
\z[{\begin{wrapfigure}[7]{r}{0.42\textwidth}%
		\vspace{-20pt}%
		\begin{flushright}%
 			\includegraphics[width=0.4\textwidth]{1.2.6}%
		\end{flushright}%
	\end{wrapfigure}%
	}]%
{При производстве варенья в большой бак постепенно наливают сироп. В первую порцию, имеющую плотность $\rho_1$, добавляют вторую, плотность которой $\rho_2$, затем третью с плотностью $\rho_3$. На графике (рис.) показано, как изменяется средняя плотность находящегося в баке сиропа по мере заполнения бака. На график капнули готовым вареньем, и часть информации пропала. Найдите массу каждой порции сиропа. До какого объёма $V_0$ был заполнен бак к тому моменту, когда средняя плотность содержимого составила $\rho_0 = 1250\,\text{кг/м}^3$?}
%
{$m_1 = 12\,\text{кг}$; $m_2 = 14\,\text{кг}$; $V_0 = 13.3\,\text{дм}^3$}

%%%%%%%%%%%%%%%%%%%%%%%%%%%%%%%%%%%%%%%%%%%%%%%%%%%%
\z{Масса шприца с $V_1 = 2\,\text{мл}$ лекарства равна $m_1 = 13.5\,\text{г}$, а его масса с $V_2 = 5\,\text{мл}$ лекарства $m_2 = 18.0\,\text{г}$. Площадь поршня шприца $S_1 = 1\,\text{см}^2$, а диаметр внутреннего отверстия иглы в $\alpha = 40$ раз меньше диаметра поршня. Определите массу пустого шприца, плотность лекарства и с какой средней скоростью лекарство выпускалось из иглы, если весь объём $V_2$ был выпущен за время $t = 20\,\text{с}$.}
%
{$u = 2\,\text{м/с}$}

